% 来源：docs/06-normalizing-flows.md
% 仅建立章节框架；正文、公式、图表与习题内容待后续迁移。
\chapter{归一化流}
\label{chap:flow}

\section*{本章要点}
% TODO：迁移原章“先抓住这几件事”的要点。
\contentplaceholder

\section{为什么需要可逆生成模型？}
\label{sec:flow-motivation}
% TODO：补充本节内容。
\contentplaceholder

\subsection{确定性双射与同维潜变量}
\label{subsec:flow-bijection}
% TODO：补充本节内容。
\contentplaceholder

\subsection{基分布、生成方向与求逆方向}
\label{subsec:flow-directions}
% TODO：补充本节内容。
\contentplaceholder

\section{变量变换：保持概率质量，修正密度}
\label{sec:flow-change-of-variables}
% TODO：补充本节内容。
\contentplaceholder

\subsection{一维例子}
\label{subsec:flow-one-dimensional}
% TODO：补充本节内容。
\contentplaceholder

\subsection{多维公式}
\label{subsec:flow-multidimensional}
% TODO：补充本节内容。
\contentplaceholder

\subsection{局部体积变化与对数密度修正}
\label{subsec:flow-volume}
% TODO：补充本节内容。
\contentplaceholder

\section{多层组合与最大似然训练}
\label{sec:flow-composition}
% TODO：补充本节内容。
\contentplaceholder

\subsection{复合变换与对数行列式求和}
\label{subsec:flow-logdet}
% TODO：补充本节内容。
\contentplaceholder

\subsection{训练、生成与潜变量推断}
\label{subsec:flow-workflows}
% TODO：补充本节内容。
\contentplaceholder

\section{如何让 Jacobian 可计算？}
\label{sec:flow-jacobian}
% TODO：补充本节内容。
\contentplaceholder

\subsection{行列式计算成本与三角结构}
\label{subsec:flow-triangular}
% TODO：补充本节内容。
\contentplaceholder

\subsection{Planar Flow 与可逆性条件}
\label{subsec:flow-planar}
% TODO：补充本节内容。
\contentplaceholder

\subsection{NICE：加性耦合与缩放}
\label{subsec:flow-nice}
% TODO：补充本节内容。
\contentplaceholder

\subsection{Real NVP：仿射耦合}
\label{subsec:flow-real-nvp}
% TODO：补充本节内容。
\contentplaceholder

\section{自回归 Flow：MAF 与 IAF}
\label{sec:flow-autoregressive}
% TODO：补充本节内容。
\contentplaceholder

\subsection{MAF：条件参数依赖历史数据}
\label{subsec:flow-maf}
% TODO：补充本节内容。
\contentplaceholder

\subsection{IAF：条件参数依赖历史噪声}
\label{subsec:flow-iaf}
% TODO：补充本节内容。
\contentplaceholder

\subsection{生成、求逆与密度评估的效率比较}
\label{subsec:flow-comparison}
% TODO：补充本节内容。
\contentplaceholder

\section{概率密度蒸馏与 Parallel WaveNet}
\label{sec:flow-distillation}
% TODO：补充本节内容。
\contentplaceholder

\subsection{自回归教师与 IAF 学生}
\label{subsec:flow-teacher-student}
% TODO：补充本节内容。
\contentplaceholder

\subsection{反向 KL 蒸馏与计算流程}
\label{subsec:flow-reverse-kl}
% TODO：补充本节内容。
\contentplaceholder

\section{能力与限制}
\label{sec:flow-strengths-limits}
% TODO：补充本节内容。
\contentplaceholder

\subsection{精确密度、可逆结构与表达能力}
\label{subsec:flow-expressivity}
% TODO：补充本节内容。
\contentplaceholder

\subsection{离散数据与连续密度的区别}
\label{subsec:flow-discrete-data}
% TODO：补充本节内容。
\contentplaceholder

\subsection{连续时间 Flow 的后续阅读}
\label{subsec:flow-future}
% TODO：补充本节内容。
\contentplaceholder

\section{自检}
\label{sec:flow-review}
% TODO：补充本节内容。
\contentplaceholder

\section{来源与延伸阅读}
\label{sec:flow-reading}
% TODO：补充本节内容。
\contentplaceholder
